\chapter{Literature Review and Background}
\label{ch:background}

This chapter reviews the concepts a reader needs before the system-specific material in Chapters~3--5 can be followed comfortably. It is deliberately written in general, textbook terms rather than in terms of the specific STM32H753 project: every mechanism introduced here is explained first as it exists in the wider embedded-systems literature, and only applied to the concrete registers, addresses, and defects of this project once Chapter~3 begins.

\section{Radar Georeferencing Principles}

Drone-based radar search-and-rescue requires translating local radar detections into global coordinates. A radar sensor reports targets in polar coordinates: slant range $r$, azimuth $\theta$, and elevation $\phi$. To map these onto a global reference frame, the sensor's own position and orientation must be known at the exact instant of measurement.

The standard approach uses an \emph{Inertial Measurement Unit (IMU)} and \emph{GPS} to determine the drone's position $\mathbf{p} = [p_n, p_e, p_d]^T$ in the North-East-Down (NED) frame and its attitude as a quaternion $\mathbf{q} = [q_w, q_x, q_y, q_z]$. The radar detection is transformed from the sensor frame to the body frame (via the sensor mounting geometry), then to the NED frame (via the attitude quaternion), and finally to the global WGS84 frame (via the GPS position).

When the drone is moving, the attitude and position must be known at the exact microsecond of the radar sweep. A timing error of 1\,ms at 3\,m/s drone velocity produces 3\,mm of spatial drift. This requirement for tight temporal synchronization motivates the use of hardware timer input capture rather than software timestamping.

\section{Kalman Filtering for Sensor Fusion}

The Extended Kalman Filter (EKF)~\cite{thrun2005,grewal2008} is the standard algorithm for fusing noisy sensor measurements with a dynamic model of system evolution. It operates in two steps:

\textbf{Prediction:} The state vector $\mathbf{x}$ is propagated forward in time using a process model $f(\mathbf{x}, \mathbf{u})$:
$$\mathbf{x}_{k|k-1} = f(\mathbf{x}_{k-1|k-1}, \mathbf{u}_k)$$
The covariance $\mathbf{P}$ is propagated using the Jacobian $\mathbf{F} = \partial f / \partial \mathbf{x}$:
$$\mathbf{P}_{k|k-1} = \mathbf{F} \mathbf{P}_{k-1|k-1} \mathbf{F}^T + \mathbf{Q}$$
where $\mathbf{Q}$ is the process noise covariance.

\textbf{Measurement update:} A measurement $\mathbf{z}$ is predicted using the measurement model $h(\mathbf{x})$:
$$\mathbf{\hat{z}} = h(\mathbf{x}_{k|k-1})$$
The innovation $\mathbf{y} = \mathbf{z} - \mathbf{\hat{z}}$ is weighted by the Kalman gain $\mathbf{K}$:
$$\mathbf{K} = \mathbf{P}_{k|k-1} \mathbf{H}^T (\mathbf{H} \mathbf{P}_{k|k-1} \mathbf{H}^T + \mathbf{R})^{-1}$$
where $\mathbf{H} = \partial h / \partial \mathbf{x}$ is the measurement Jacobian and $\mathbf{R}$ is the measurement noise covariance. The state and covariance are then updated:
$$\mathbf{x}_{k|k} = \mathbf{x}_{k|k-1} + \mathbf{K} \mathbf{y}$$
$$\mathbf{P}_{k|k} = (\mathbf{I} - \mathbf{K} \mathbf{H}) \mathbf{P}_{k|k-1}$$

For radar-inertial georeferencing, the state vector is typically 6--15 dimensions (position, velocity, attitude, and optionally sensor biases). This thesis uses a simplified 6-state model ($[p_n, p_e, p_d, v_n, v_e, v_d]^T$) with attitude fused separately via quaternion, reducing computational requirements while maintaining the core fusion architecture.

\section{Quaternion-Based Coordinate Transforms}

Rotations in 3D space can be represented by Euler angles, rotation matrices, or quaternions. Quaternions~\cite{kuipers1999} are preferred for embedded systems because they avoid gimbal lock, require only 4 floats (vs.\ 9 for a matrix), and compose efficiently.

A quaternion $\mathbf{q} = [q_w, q_x, q_y, q_z]$ represents a rotation of angle $\theta$ about unit axis $\hat{\mathbf{n}} = [n_x, n_y, n_z]$:
$$q_w = \cos(\theta/2), \quad q_x = n_x \sin(\theta/2), \quad q_y = n_y \sin(\theta/2), \quad q_z = n_z \sin(\theta/2)$$

A vector $\mathbf{v}$ is rotated by $\mathbf{q}$ using $\mathbf{v}' = \mathbf{q} \mathbf{v} \mathbf{q}^{-1}$ where $\mathbf{q}^{-1} = [q_w, -q_x, -q_y, -q_z]$ (the conjugate) for unit quaternions. A common defect is applying the conjugate when the forward quaternion is needed (or vice versa), which rotates in the opposite direction.

The coordinate transform pipeline for radar georeferencing is:
\begin{enumerate}
\item \textbf{Polar to Cartesian:} $(r, \theta, \phi) \rightarrow (x, y, z)$ in the radar frame
\item \textbf{Radar to body:} Apply sensor mounting rotation
\item \textbf{Body to NED:} Rotate by attitude quaternion $\mathbf{q}_{\text{attitude}}$
\item \textbf{NED to WGS84:} Apply reference GPS position
\end{enumerate}

\section{Telemetry Framing and I/O Strategies}

A microcontroller that receives data from a serial peripheral, such as a MAVLink-family~\cite{mavlink} telemetry stream, has three strategies available. The simplest is \emph{polling}: software repeatedly reads a status register until a ``data ready'' bit becomes set, then reads the data register. The second is \emph{interrupt-driven I/O}: the peripheral raises a hardware interrupt when data arrives. The third, \emph{Direct Memory Access (DMA)}, removes the CPU from the transfer path entirely: a dedicated hardware block moves bytes from the peripheral's data register directly into memory, notifying the CPU only when a complete transfer has finished.

All three strategies rest on memory-mapped I/O (MMIO): a peripheral register is an address wired directly to physical logic rather than to a memory cell. Reading it can have side effects (clearing a flag, releasing a hardware buffer slot); writing to it can cause a real-world action (driving a pin, starting a transfer).

\section{The Arm Cortex-M7 Core and Exception Model}

The Arm Cortex-M7~\cite{armcm7trm} is a high-performance member of Arm's Cortex-M family: a six-stage, superscalar pipeline implementing the Thumb-2 instruction set, with optional L1 instruction and data caches. Two details matter throughout this thesis. First, the processor maintains two stack pointers: the Main Stack Pointer (MSP) and the Process Stack Pointer (PSP). Second, the processor has two execution modes: \emph{Thread Mode} for ordinary application code and \emph{Handler Mode} for exception/interrupt handlers.

Every Cortex-M exception is serviced by fetching a handler address from the \emph{vector table}, indexed by exception number. The first sixteen entries are processor architectural exceptions; entries from index 16 onward correspond to device-specific interrupt requests routed through the \emph{Nested Vectored Interrupt Controller (NVIC)}. The NVIC arbitrates among interrupt sources, comparing priorities and supporting \emph{tail-chaining} for low-latency interrupt handling.

Before any of this machinery can run, the processor must boot. Every Cortex-M device follows the same fixed sequence: the core loads its initial stack pointer from the first word of the vector table, then loads the second table word into the program counter to begin executing the reset handler, which copies pre-initialised globals from Flash to RAM, zeroes uninitialised globals, and calls \texttt{main}.

\section{The STM32H7 Peripheral Subsystem}

STMicroelectronics' STM32H7 family organises its on-chip peripherals as memory-mapped register blocks whose clock signal is not enabled by default. Software must explicitly enable a peripheral's clock through the Reset and Clock Control (RCC) register before that peripheral will respond to any configuration.

Three specific peripheral subsystems recur throughout this thesis: a \emph{Universal Synchronous/Asynchronous Receiver-Transmitter (USART)} for the MAVLink telemetry link; a 32-bit general-purpose \emph{timer} (TIM2) for microsecond-precision frame-sync timestamping; a hardware \emph{Cyclic Redundancy Check (CRC)} accelerator for packet integrity verification; a \emph{Serial Peripheral Interface (SPI)} for output streaming; and a \emph{Memory Protection Unit (MPU)} for memory safety. Each is a well-understood peripheral category; what makes each a source of genuine defects is a specific, documented behaviour of STMicroelectronics' implementation.

\section{Hardware Emulation and Renode}

Testing embedded firmware exclusively on physical hardware is slow, hard to automate, and destructive when a defect is capable of corrupting memory. Hardware emulation executes the same compiled binary inside a software model of the chip's processor core, memory map, and peripherals.

Renode~\cite{renode}, developed by Antmicro and released under the MIT licence, is a \emph{functional} simulator: it executes real target machine code against modelled peripherals with enough fidelity to run production firmware unmodified, while explicitly not attempting cycle-accurate timing. Renode provides full determinism of execution and a shared virtual-time model, making it attractive for reproducible research.

The property that matters most for this thesis is the direct consequence of Renode's position in the simulator design space: because it is a functional, not cycle-accurate, model, there exist hardware behaviours it does not reproduce with complete fidelity --- a cache's eviction timing, a peripheral's under-modelled internal state machine, a bus arbiter's exact contention behaviour. This thesis documents these divergences explicitly and engineers the firmware defensively for physical silicon.

\section{Common Classes of Embedded-C and ABI Defects}

Three defect classes recur directly in Chapter~4:

\textbf{Structure padding.} A C compiler inserts unused padding bytes between structure members to satisfy alignment requirements. This is invisible for structures used only in memory, but hazardous for structures whose byte layout is also defined by an external wire protocol. The compiler's packed attribute suppresses this behaviour.

\textbf{Cache coherency with a second bus master.} When a DMA controller writes to a memory location the CPU has already cached, the CPU's cached copy becomes stale. The CPU must explicitly invalidate it via a cache-maintenance instruction before its next read.

\textbf{Endianness and byte-order convention.} When a multi-byte value crosses a boundary between two components that do not agree on byte order (e.g., a little-endian ARM core and a hardware accelerator with its own internal convention), the result is silently wrong. This must be checked against each specific peripheral's documented behaviour.